% Generated draft: requires mathematical and source review.
\documentclass[11pt]{article}
\usepackage[margin=25mm]{geometry}
\usepackage{fontspec}
\setmainfont{Latin Modern Roman}
\usepackage{amsmath,amssymb,mathtools,unicode-math}
\setmathfont{Latin Modern Math}
\usepackage{enumitem,needspace,xurl,hyperref}
\setlength{\parindent}{0pt}
\setlength{\parskip}{5pt}
\setlength{\emergencystretch}{3em}
\newcommand{\N}{\mathbb{N}}
\newcommand{\Z}{\mathbb{Z}}
\newcommand{\Q}{\mathbb{Q}}
\newcommand{\R}{\mathbb{R}}
\newcommand{\C}{\mathbb{C}}
\newcommand{\F}{\mathbb{F}}
\newcommand{\A}{\mathbb{A}}
\newcommand{\PP}{\mathbb{P}}
\newcommand{\E}{\mathbb{E}}
\newcommand{\eps}{\varepsilon}
\newcommand{\dd}{\,\mathrm{d}}
\begin{document}
\section*{UnsolvedMath ID: 800332}
\textbf{Source problem code:} catalogue entry
\subsection{Definitions and mathematical statement}
A unit lozenge is the union of two unit equilateral triangles sharing an edge in the triangular lattice. A lozenge tiling of a lattice region is a covering by unit lozenges with pairwise disjoint interiors.

Under the conventional triangular-lattice interpretation, let \(H_n\), whenever defined, be an equiangular hexagon with consecutive side lengths
\[
2n,\quad 2n+3,\quad 2n,\quad 2n+3,\quad 2n,\quad 2n+3,
\]
measured in lattice units. Its three long sides have length \(2n+3\). Write \(T_{n,1},T_{n,2},T_{n,3}\) for the triangular portions to be removed from the middles of those sides, and write
\[
R_n=H_n\setminus\bigcup_{i=1}^{3}T_{n,i}.
\]
The boundary convention in this notation is immaterial to tiling interiors; the intended remaining lattice cells must still be specified.

\textbf{Enumeration problem attributed to Propp (1998), Problem 3.} Determine the number
\[
M(R_n)=\#\{\text{lozenge tilings of }R_n\}
\]
for every intended admissible value of \(n\).

\textbf{Unresolved qualifications.} The supplied source excerpt does not specify the sizes or precise lattice placement of the three removed triangles, whether their sizes agree, or the intended meaning of their being in the middle of a side. Consequently, the notation \(T_{n,i}\) represents unresolved geometric data, not a uniquely defined construction. Equilateral lattice triangles are a conventional interpretation of the removals, but that interpretation is not explicitly stated in the excerpt. The domain of \(n\) is also unstated; taking \(n\) to be a positive integer would be an additional indexing assumption, and degenerate cases are not addressed. Thus the excerpt does not determine a unique sequence \(M(R_n)\).
\subsection{Short English statement}
Propp's supplied Problem 3 asks for the number of lozenge tilings of a hexagon with alternating side lengths \(2n\) and \(2n+3\), after removing a triangle from the middle of each long side. The triangle sizes and placements, and the domain of \(n\), require clarification before the region and its tiling count are uniquely specified.
\subsection{Sources}
\begin{itemize}
\item Local UnsolvedMath registry, AMR Open Problem Lists: permanent ID \(800332\), exact source problem code \texttt{catalogue entry}, title \emph{Enumeration of Matchings --- Problem 3}. Supplied upstream revision: \texttt{d371ee917ac9647fbb06cd6fc529deec6f3dda82}.
\item Propp, \emph{Twenty Open Problems in Enumeration of Matchings} (1998), Problem 3 (\#3). Supplied source URL: \url{https://arxiv.org/abs/math/9801060}. The registry attributes the problem to Propp and gives proposed year 1998.
\item The supplied provenance records access on 2026-10-16 and extraction by \texttt{arxiv\_latex}. This is supplied provenance, not an independently verified source visit. The registry record, created and updated on 2026-10-16, labels the problem \texttt{open}; this supplied status is not an independently verified mathematical or literature assessment.
\end{itemize}
\end{document}
